\documentclass[11pt,a4paper]{article}
\usepackage[utf8]{inputenc}
\usepackage[T1]{fontenc}
\usepackage[portuguese]{babel}
\usepackage{xcolor}
\usepackage{hyperref}
\usepackage{geometry}
\usepackage{tcolorbox}

\geometry{margin=2.5cm}

\definecolor{primary}{RGB}{39, 174, 96}
\definecolor{secondary}{RGB}{44, 62, 80}
\definecolor{codebg}{RGB}{240, 240, 240}

\hypersetup{colorlinks=true, linkcolor=primary, urlcolor=primary}

\title{\color{secondary}\Huge\textbf{Solução: Exercício 5.4 - Wikipedia com init e sidecar}}
\author{\color{primary}DevOps com Kubernetes -- 2026}
\date{}

\begin{document}
\maketitle

\section*{\color{primary}Guia de Solução}
\begin{tcolorbox}[colback=green!5!white,colframe=green!60!black,title=Resolução Passo a Passo]
### Passo a Passo

1. \textbf{Criar Volume}: Crie um manifesto de Deployment com um volume do tipo \texttt{emptyDir} para compartilhar os dados entre os contêineres do Pod. Ele será montado em \texttt{/usr/share/nginx/html}.
2. \textbf{Init Container}: Na secção \texttt{initContainers}, configure um contêiner baseado em \texttt{curlimages/curl}. O comando deverá ser: \texttt{sh -c "curl -L -o /usr/share/nginx/html/index.html https://en.wikipedia.org/wiki/Kubernetes"}.
3. \textbf{Sidecar Container}: Na seção regular \texttt{containers}, crie um contêiner secundário usando \texttt{curlimages/curl}. Em vez de morrer, ele deve rodar um loop longo: \texttt{sh -c "while true; do sleep \$((RANDOM \% 600 + 300)); curl -L -o /usr/share/nginx/html/index.html https://en.wikipedia.org/wiki/Special:Random; done"}.
4. \textbf{Main Container}: Crie o contêiner final usando \texttt{nginx:alpine} com o \texttt{volumeMounts} devidamente exposto, que apenas servirá os arquivos HTML alterados de maneira contínua.
\end{tcolorbox}

\end{document}
